\section{Introducción}
La Unidad Aritmético-Lógica (ALU, por sus siglas en inglés Arithmetic Logic Unit) es uno de los componentes fundamentales de la arquitectura de cualquier procesador digital. Se trata del bloque combinacional (o secuencial, según el diseño) encargado de ejecutar las operaciones aritméticas y lógicas que constituyen la base del procesamiento de datos en un sistema digital, formando parte esencial de la unidad central de procesamiento (CPU) junto con la unidad de control y los registros.

Entre las operaciones típicas que realiza una ALU se encuentran las aritméticas —como la suma, la resta, la multiplicación y, en algunos casos, la división— y las lógicas, entre las que se destacan las operaciones AND, OR, XOR, NOT, así como operaciones de desplazamiento (shift) y rotación de bits. La selección de la operación a ejecutar se realiza generalmente mediante una señal de control u opcode, que indica a la ALU qué función debe aplicar sobre los operandos de entrada en un momento dado.

Desde el punto de vista del diseño digital, la ALU representa un caso de estudio particularmente relevante, ya que combina conceptos de lógica combinacional, diseño modular y optimización de recursos. Su implementación mediante lenguajes de descripción de hardware (HDL), como Verilog, permite modelar su comportamiento a distintos niveles de abstracción —comportamental, de flujo de datos o estructural— y posteriormente sintetizarla para su implementación física en dispositivos programables (FPGA) o circuitos integrados de aplicación específica (ASIC).

En este informe se describe el diseño de una ALU en Verilog y su implementación en la placa Basys 3. Primero se presenta la especificación planteada por el enunciado y luego el diseño de los módulos que componen la solución. A continuación se describe la verificación mediante simulación y, por último, se analizan los resultados de la síntesis e implementación en cuanto a utilización de recursos y tiempos.
